
\begin{abstract}
本文提出一种离散且可计算的“折叠--回溯”闭包图模型：将长度为 $m$ 的二进制微观态集合 $X_m=\{0,1\}^{m}$ 通过 Zeckendorf 表示的\textbf{截断投影}映射到稳定宏观态集合 $Y_m$（无相邻 $1$ 的 Fibonacci/Zeckendorf 字），得到 $2^{m}\to F_{m+2}$ 的信息压缩；被截断的高位 Zeckendorf 尾部作为“时间纤维”保留。我们把折叠边与微观邻接（Hilbert-lattice、bit-split lattice、hypercube 等）合并成带颜色的闭包图，并把其图同构问题等价翻译到 Babai 体系中的 \textbf{\WL{}--相干配置}语言。通过穷举实验（$m=3$ 至 $18$）发现：一维/二维薄格导致 \WL{} 塌缩显著变慢，而三维及更高“适度维数”的格使塌缩步数近似常数；同时在少数“共振” $m$ 上出现残余 $\ZZ_2$ 纤维对称（不可恢复的正交时间位）。本文给出 \GI{}/\WL{} 的严格翻译定理、残余对称的可验证判据、以及可复现实验指标与开放问题，为将 $64\to 21$（$m=6$）作为最小闭合版本的离散内核提供可发表的数学与计算框架。
\end{abstract}

\begin{sloppypar}
\textbf{关键词：} 图同构；Weisfeiler--Leman；相干配置；Zeckendorf 表示；Fibonacci 编码；Hilbert-lattice；闭包图；残余对称；维数相变
\end{sloppypar}

\paragraph{约定与记号。}
除非另行说明，$\log$ 表示自然对数。长度为 $m$ 的二进制串写作 $x=x_0x_1\cdots x_{m-1}$，其中 $x_i\in\{0,1\}$。集合的基数记为 $\abs{\cdot}$。本文中的 Fibonacci 数列采用 Zeckendorf 版本：$F_1=1$，$F_2=2$，$F_{k+2}=F_{k+1}+F_k$。关于 \GI{}、\WL{} 与相干配置的背景可参见 \cite{babai2016gi,wl2018symmetryvsregularity}；Zeckendorf 表示见 \cite{zeckendorf1972representation}。

